% !PRINTED 24
\weberpage{24}
%<<CONT
dass die Indexreihen $j_{r+1}, j_{r+2} \ldots j_\mu$ und $j_{r+1}',
j_{r+2}' \ldots j_\mu'$, von der Ordnung abgesehen, übereinstimmen.
Folglich müssen auch die vorangehenden Theile der Indexreihen
\[
j_1, j_2 \ldots j_r \quad \text{und} \quad j_1', j_2' \ldots j_r'
\]
übereinstimmen.

\begin{center}
\S.\,7.

\gesperrt{Weitere S\"atze \"uber die Compositionsreihen.}
\end{center}

Von den Compositionsreihen gelten noch mannigfaltige Sätze, von
denen wir hier noch die wichtigsten ableiten.

\gesperrt{II. Ist $P_\nu$ irgend ein Normaltheiler von $P$, so giebt
es eine Compositionsreihe von $P$, in der $P_\nu$ vorkommt.}

Ist $P_\nu$ ein grösster Normaltheiler von $P$, so können wir $P_\nu =
P_1$ setzen und die Compositionsreihe von da an beliebig
fortsetzen. Ist aber $P_\nu$ kein grösster Normaltheiler von $P$, so
suchen wir einen grössten Normaltheiler über $P_\nu$, den wir mit
$P_1$ bezeichnen. Dann ist $P_\nu$ ein Normaltheiler von $P_1$. Ist
es ein grösster, so setzen wir $P_\nu = P_2$ und setzen von da die
Compositionsreihe beliebig fort. Ist aber $P_\nu$ noch kein grösster
Normaltheiler von $P_1$, so suchen wir einen grössten
Normaltheiler von $P_1$ über $P_\nu$ und fahren so fort, bis wir
schliesslich zu $P_\nu$ gelangen, was nach einer endlichen Anzahl von
Schritten nothwendig eintreten muss. Wenn wir dann eine
Compositionsreihe von $P_\nu$ anhängen, so haben wir eine
Compositionsreihe von $P$, in der $P_\nu$ vorkommt, wie es der Satz
II. verlangt.

In einer Compositionsreihe einer Gruppe $P$ ist, vermöge der
Definition, jedes Glied ein Normaltheiler jedes vorangehenden. Das
erste Glied $P_1$ und das letzte Glied 1 sind zugleich
Normaltheiler von $P$ selbst. Bei den zwischenliegenden $P_2, P_3
\ldots$ wird dies im Allgemeinen nicht der Fall sein. In der
Compositionsreihe giebt es also gewisse ausgezeichnete Glieder, die
zugleich Normaltheiler von $P$ selbst sind, die eine besondere
Rolle spielen\footnote{Sie bilden die sogenannte Hauptreihe von $P$
(\gesperrt{Netto}).}. Hierüber leiten wir im Folgenden noch einen
wichtigen Satz ab.

Angenommen, es sei $Q$ ein Glied einer Compositionsreihe von $P$,
das zugleich Normaltheiler von $P$ ist, während das auf $Q$
%CONT>>
